\section{Design Specifications}\label{sec:design_specifications}

\subsection{System Architecture}
Figure~\ref{fig:overall_architecture} shows the complete design. The camera and
wearables use separate input paths because their rates and errors are different.
The vision path creates pitch tracks and identity evidence. The sensor path
creates timestamped wearable readings with quality states. Clock alignment and
player binding join them on the project PC.

\begin{figure}[H]
\centering
\includegraphics[width=0.96\textwidth,height=0.69\textheight,keepaspectratio]{02_overall_architecture.png}
\caption{Overall system architecture}
\label{fig:overall_architecture}
\end{figure}

The joined event store in Figure~\ref{fig:overall_architecture} keeps raw source
references, mapped time, player identity, and quality. Analytics reads these
events and gives results to FastAPI. The local React dashboard receives live
tracks, annotated frames, sensor values, and tactical updates. Reviewed exports
go to the training pipeline, and a trained model is used only after its test
result and version are recorded.

The VPS relay is part of the deployed architecture. The wearable sends raw
timestamped batches to \texttt{cpg44.nivaspms.com}. Cloudflare and Caddy route
them to a separate relay process. The relay validates each event, adds receive
time and a relay sequence, then provides live and replay endpoints to the local
sensor processor. Its default cache is 4 MiB in memory and it has no database
or file volume.

\subsubsection{Wearable Component}
The field unit uses ESP32-S3 GPIO 4 and 5 for the I2C sensor bus. The NEO-6M
transmit line connects to receive GPIO 16, with GPIO 17 available for the
reverse serial direction. The firmware sends device ID, boot ID, sequence,
device time, sensor values, and status. Wi-Fi details and player tags are set
through provisioning and are not fixed in the source code.

The host keeps raw acceleration and a filtered magnitude. GPS includes fix and
age information. PPG output includes contact, motion, window age, signal
quality, and estimator version. When a quality test fails, the raw evidence is
kept but the derived BPM or SpO$_2$ is not treated as valid.

\subsubsection{Vision and Identity Component}
The detector uses the four football classes in the training data. ByteTrack
creates short-term tracks. Pitch points, team-colour samples, appearance
features, jersey votes, and operator corrections are added after tracking. A
track can stay unresolved when the evidence conflicts.

The ball uses a shorter recovery window and field-speed checks. This prevents a
long guessed path during occlusion. Identity corrections are saved with their
time and reason, so they can also become reviewed training labels.

\subsubsection{Relay and Reconnection Component}
The wearable publishes batches to the relay with a stable event ID, source
sequence, source time, match ID, player ID, sample type, and bounded tags. The
relay checks the token, timestamp, value range, body size, and source type.
Duplicate event IDs keep their earlier relay sequence.

The local sensor processor reconnects with its last accepted relay sequence.
Events still in the 4 MiB cache are replayed in order. If the requested sequence
has left the cache, the relay marks a cache gap and the processor starts from
the latest available state. This supports short network loss without turning
the VPS into session storage.

\subsubsection{Analytics and Training Component}
Every feature has a definition version and test inputs. Team width, length,
centroid, compactness, line spacing, pressure distance, player speed, distance,
and acceleration load are calculated from checked observations. An event points
back to the positions and feature values that caused it.

Workload model data uses one row per player-session. The dataset keeps session
date, player group, label source, missing values, and split. A simple baseline is
tested before a larger model. Injury-related output remains a workload concern
because load ratios alone should not be used as a fixed injury boundary
\cite{bowen2020acwr,dalen2021acwr}.

\subsection{Design Level Diagrams}

\subsubsection{Overall Use-Case Diagram}
Figure~\ref{fig:overall_use_case} gives the main users and actions. The coach
prepares the roster, watches the session, reviews events, and confirms
corrections. The analyst calibrates the camera, pairs devices, records data, and
exports reviewed labels. The administrator checks streams and model runs.

\subsubsection{Live-Session Activity Diagram}
Figure~\ref{fig:live_activity} shows the main live case. Camera, roster, field
mapping, and wearable checks happen before the clock starts. Video and sensor
processing then run together. Their outputs are joined before the dashboard and
event timeline are updated.

\begin{figure}[H]
\centering
\begin{subfigure}[t]{0.43\textwidth}
\centering
\includegraphics[width=\linewidth,height=0.48\textheight,keepaspectratio]{01_overall_use_case.png}
\caption{Overall use cases}
\label{fig:overall_use_case}
\end{subfigure}
\hfill
\begin{subfigure}[t]{0.53\textwidth}
\centering
\includegraphics[width=\linewidth,height=0.48\textheight,keepaspectratio]{03_live_session_activity.png}
\caption{Live-session activity}
\label{fig:live_activity}
\end{subfigure}
\caption{Overall user actions and the main live-session flow}
\end{figure}

The dependency arrows in Figure~\ref{fig:overall_use_case} show that a fused
recording needs calibration and device pairing. Training export also needs a
review step. A camera track can still exist without a wearable, but no sensor
value is assigned to that player.

At the end of Figure~\ref{fig:live_activity}, corrections and checked data are
saved before session close. Fixed queue limits prevent a slow consumer or short
network loss from increasing memory without a limit.

\subsubsection{Synchronisation Sequence Diagram}
Figure~\ref{fig:sync_sequence} explains the clock and fusion steps. SNTP gives
the wearable a UTC anchor before streaming. Each event keeps its source time,
relay receive time, and sequence. Relay delay and sequence gaps remain visible
in the processed record.

\begin{figure}[H]
\centering
\includegraphics[width=0.96\textwidth,height=0.63\textheight,keepaspectratio]{04_sync_sequence.png}
\caption{Sequence diagram for wearable and vision synchronisation}
\label{fig:sync_sequence}
\end{figure}

The fusion step in Figure~\ref{fig:sync_sequence} checks time error, identity,
packet age, and signal quality. If a check fails, the dashboard receives a
partial update with a reason. It does not receive a made-up normal value.

\subsubsection{Domain and Data Model}
The domain model in Figure~\ref{fig:domain_model} keeps observations,
corrections, features, and events separate. Device bindings have a time range
because the roster or pairing can change. A track may resolve to one player or
remain unresolved.

\begin{figure}[H]
\centering
\includegraphics[width=0.95\textwidth,height=0.62\textheight,keepaspectratio]{05_domain_model.png}
\caption{Core domain and data model}
\label{fig:domain_model}
\end{figure}

Derived features in Figure~\ref{fig:domain_model} point to their observations
and definition version. A tactical event points to its supporting features.
This allows the team to recalculate a result after a mapping or identity
correction.

The compact live envelope is shown in Table~\ref{tab:event_envelope}. Video
frames are delivered by the local API; the VPS relay accepts only validated raw
wearable events.

\begin{table}[H]
\centering
\caption{Common fields in a live event}
\label{tab:event_envelope}
\small
\begin{tabularx}{\textwidth}{p{1.4in}p{1.1in}X}
\toprule
\textbf{Field} & \textbf{Type} & \textbf{Purpose}\\
\midrule
session\_id & string & Separates recording sessions\\
source and entity & string & Identifies camera, wearable, track, player, or event\\
source\_time & integer & Time created at the source\\
session\_time & integer & Source time mapped to the session clock\\
receive\_time & integer & Host or relay arrival time\\
sequence & integer & Finds gaps, duplicates, and replay order\\
quality & object & Validity, confidence, age, timing error, and reason\\
payload & object & Typed values for the event\\
\bottomrule
\end{tabularx}
\end{table}

\subsubsection{Overall and Identity State Diagrams}
Figure~\ref{fig:session_state} gives the state of the complete session. A draft
must pass camera, roster, and calibration checks before recording. A stale
important stream changes the state to degraded, while valid streams continue.
Closing flushes bounded queues before review.

Player identity is a significant object because all joined statistics depend
on it. Figure~\ref{fig:identity_state} starts at unresolved. Team evidence gives
only a team state. Jersey or appearance evidence gives a candidate. Strong
votes or analyst review are needed for a confirmed identity.

\begin{figure}[H]
\centering
\begin{subfigure}[t]{0.43\textwidth}
\centering
\includegraphics[width=\linewidth,height=0.43\textheight,keepaspectratio]{06_session_state.png}
\caption{Overall session state}
\label{fig:session_state}
\end{subfigure}
\hfill
\begin{subfigure}[t]{0.53\textwidth}
\centering
\includegraphics[width=\linewidth,height=0.43\textheight,keepaspectratio]{07_identity_state.png}
\caption{Track-to-player identity state}
\label{fig:identity_state}
\end{subfigure}
\caption{State diagrams for the overall system and a significant identity object}
\end{figure}

Conflicting evidence in Figure~\ref{fig:identity_state} changes the state to
ambiguous. A short lost state allows recovery after an occlusion. When the
recovery limit expires, the new track starts unresolved instead of carrying an
old player name.

\subsubsection{Deployment Diagram}
Figure~\ref{fig:deployment} shows the physical boundary. The camera is connected
to the project PC, while the wearable sends through Wi-Fi to the required domain
route. Cloudflare and Caddy pass the raw sensor stream to the VPS relay. The
project PC reads this stream and runs processing, local files, API, and the
complete coaching dashboard.

\begin{figure}[H]
\centering
\includegraphics[width=0.94\textwidth,height=0.61\textheight,keepaspectratio]{08_deployment.png}
\caption{Deployment with the required memory-only VPS relay}
\label{fig:deployment}
\end{figure}

The wearable and local sensor processor in Figure~\ref{fig:deployment} use HTTPS
or WSS at \texttt{cpg44.nivaspms.com}, never the VPS IP. A token is checked
before data is accepted or returned. A restart clears the replay cache, while
full session evidence stays on the project PC.

\subsection{User Interface Diagrams}
The interface uses a light background, dark text, muted borders, and one small
accent colour. It has no neon colour, glass effect, or dark default. Status is
written as ``Live'', ``Delayed'', ``Low signal'', or ``Disconnected'', so colour
is not the only meaning.

The main pages are Live Session, Session Review, Calibration and Identity,
Wearables, and Models and Data. The live pitch stays largest. Cards move below
the pitch on a narrow screen. Technical model controls are kept away from the
coach's live view.

\evidenceplaceholder{Live dashboard with pitch markings and player data}{Insert the approved live-session screen capture here.}{dashboard_live}

Figure~\ref{fig:dashboard_live} will show the large pitch, player and ball
markings, session clock, stream health, and selected-player panel. Player labels
must remain readable on the panel projector. A weak identity is shown as
unresolved.

\evidenceplaceholder{Session analytics and tactical timeline}{Insert the approved session-review screen capture here.}{dashboard_analytics}

Figure~\ref{fig:dashboard_analytics} will show a selected event, the matching
pitch state, team geometry, and player workload. Only the values required to
explain the event are shown first; more detail is opened when needed.

\evidenceplaceholder{Wearable quality and live sensor page}{Insert the approved wearable screen capture here.}{dashboard_wearable}

Figure~\ref{fig:dashboard_wearable} will show device-player binding, packet age,
sequence gaps, contact, motion quality, GPS fix, BPM, SpO$_2$, and acceleration.
Quality stays beside the measurement on both wide and narrow layouts.

\subsection{Snapshots of Working Prototype}
The panel demonstration follows one fixed sequence:

\begin{enumerate}
  \item Show the wearable, sensor placement, and safe fitting.
  \item Start the camera, load the roster, and verify pitch calibration.
  \item Pair the wearable to a visible player and start the session clock.
  \item Show live player and ball markings, telemetry, and stream health.
  \item Show one aligned event, one identity correction, and its tactical and workload evidence.
  \item Reconnect the local sensor processor, close the session, and show the checked export.
\end{enumerate}

\begin{figure}[H]
\centering
\includegraphics[width=0.88\textwidth,height=0.48\textheight,keepaspectratio]{training_results.png}
\caption{Training curves from the recorded 30-epoch detector run}
\label{fig:training_results}
\end{figure}

Figure~\ref{fig:training_results} is evidence that the detector training path
runs. Its final recorded values are precision 0.8435, recall 0.7318,
AP$_{50}$ 0.7476, and AP$_{50:95}$ 0.4589 on that split. These are not claimed as
final campus-field accuracy.

\evidenceplaceholder{Assembled ESP32-S3 football wearable}{Insert the labelled hardware prototype photograph here.}{hardware_front}

Figure~\ref{fig:hardware_front} will show the enclosure, strap, optical contact
surface, power control, and cable strain relief. Labels should make the sensor
connections easy to explain.
